% =====================================================================
%  LAMPIRAN B — PENGGUNAAN BAHASA
%  Isi diambil langsung dari Panduan Skripsi Filkom UB v3.0 (wajib
%  dilampirkan verbatim di setiap skripsi/proposal, bukan hasil tulisan
%  penulis). Jangan diubah kecuali panduan resminya sendiri berubah.
% =====================================================================

\chapter*{Lampiran B Penggunaan Bahasa}
\addcontentsline{toc}{chapter}{Lampiran B Penggunaan Bahasa}

Bahasa yang dipakai dalam naskah adalah Bahasa Indonesia yang baku.
Setiap kalimat berita harus memiliki subjek dan predikat, dan umumnya
dilengkapi objek, pelengkap, atau keterangan. Setiap paragraf biasanya
terdiri dari beberapa kalimat. Penuturan isi dalam kalimat, paragraf,
maupun antarparagraf harus menggunakan bahasa yang tepat dan
menggambarkan alur logika yang runtut.

Penulisan bahasa asing yang sudah diserap dalam Bahasa Indonesia
disesuaikan dengan kaidah Bahasa Indonesia. Sedapat mungkin dihindari
penggunaan istilah asing jika padanan dalam Bahasa Indonesia sudah ada.
Jika terpaksa menggunakan istilah asing, penulisannya harus sesuai ejaan
aslinya dan dicetak miring (\emph{italic}), kecuali istilah tersebut
adalah nama.

Sebagai referensi penulisan Bahasa Indonesia yang baku, dokumen berikut
dapat digunakan:

\begin{itemize}
  \item Kamus Bahasa Indonesia, Tim Penyusun, Pusat Bahasa Departemen
        Pendidikan Nasional, Jakarta, 2008.
  \item Peraturan Menteri Pendidikan Nasional Republik Indonesia Nomor
        46 Tahun 2009 tentang Pedoman Umum Ejaan Bahasa Indonesia yang
        Disempurnakan.
  \item Kamus Besar Bahasa Indonesia dalam jaringan (KBBI daring):
        \url{https://kbbi.web.id}
\end{itemize}
